% =============================================================================
% HAUPTDOKUMENT (main.tex)
% =============================================================================

% GEÄNDERT: 'footsepline=0.4pt' hinzugefügt
\documentclass[11pt, a4paper, headsepline=0.4pt, footsepline=0.4pt]{scrartcl}

% --- EINSTELLUNGEN LADEN ---
% =============================================================================
% PRÄAMBEL (praeambel.tex)
% -----------------------------------------------------------------------------
% Lädt alle Pakete und globalen Einstellungen.
% =============================================================================

% --- GRUNDEINSTELLUNGEN ---
\usepackage[utf8]{inputenc}
\usepackage[T1]{fontenc}
\usepackage[ngerman]{babel}
\usepackage{graphicx}
\usepackage{csquotes}

% --- LITERATURVERZEICHNIS (BIBLATEX) ---
\usepackage[backend=biber, style=alphabetic, sorting=ynt]{biblatex}

% --- MATHEMATIK ---
\usepackage{amsmath, amssymb, amsthm, latexsym, derivative}

% --- GRAFIKEN & DIAGRAMME ---
\usepackage{epsfig, tikz, subcaption, caption, adjustbox, eso-pic}
\usepackage[european]{circuitikz}
\usepackage{siunitx}
\sisetup{locale = DE}

\usepackage{svg}
\usepackage{filecontents}

% --- TABELLEN ---
\usepackage[table]{xcolor} % table-Option für colortbl-Kompatibilität
\usepackage{booktabs, multirow, colortbl}

% --- SONSTIGE PAKETE ---
\usepackage{enumitem}
\usepackage{makeidx}
% \makeindex % Vorerst auskommentiert
\usepackage{totpages} % Robustes Paket für Gesamtseitenzahl
\usepackage[colorlinks=true, linkcolor=blue, citecolor=green, urlcolor=magenta, hidelinks]{hyperref}
\usepackage[acronym, nonumberlist, toc, section]{glossaries}
\makeglossaries
\usepackage{chngcntr}
\usepackage[normalem]{ulem}
\usepackage{calc}
\usepackage{pdfpages}
\usepackage{blindtext}
\usepackage{float}
\usepackage{newfloat}
\usepackage[most]{tcolorbox}
\usepackage[framemethod=TikZ]{mdframed}
\usepackage{etoolbox}

% --- LAYOUT-DATEI LADEN ---


% --- BENUTZERDEFINIERTE EINSTELLUNGEN ---
\setlength{\parindent}{0em}
\setlength{\parskip}{0.5em}

% --- DEINE FLOAT-UMGEBUNGEN ---
\DeclareFloatingEnvironment[listname={Diagrammverzeichnis}, name=Diagramm, fileext=lod, placement=htp]{diagram}
\DeclareFloatingEnvironment[listname={Oszillogrammverzeichnis}, name=Oszillogramm, fileext=loo, placement=htp]{oszillo}

% --- ZÄHLER-EINSTELLUNGEN ---
\counterwithin{figure}{section}
\counterwithin{table}{section}
\counterwithin{oszillo}{section}
\counterwithin{diagram}{section}

% --- DEINE BOX-DEFINITIONEN ---
\newmdenv[linecolor=red, linewidth=1pt, frametitle=Hinweis, frametitlebackgroundcolor=red, frametitlefont=\color{white}\bfseries, frametitlerule=true, backgroundcolor=white, roundcorner=2pt]{hinweisbox}
\newmdenv[linecolor=gray, linewidth=1pt, frametitle=Aufgabe, frametitlebackgroundcolor=white, frametitlefont=\color{black}\bfseries, frametitlerule=true, backgroundcolor=white, roundcorner=2pt]{aufgabenbox}
\newcounter{rechnungbox}[section]
\renewcommand{\therechnungbox}{\arabic{rechnungbox}}
\newtcolorbox{rechnungbox}[1][]{colframe=darkgray, colback=white, coltitle=white, colbacktitle=gray, title={Rechnung \therechnungbox}, fonttitle=\bfseries, boxed title style={size=small},
    boxrule=1pt, sharp corners=south, before upper=\stepcounter{rechnungbox}, #1}

% --- DEINE TOGGLES ---
\newtoggle{showtoc}
\newtoggle{showtables}
\newtoggle{showfigures}
\newtoggle{showoszillos}
\newtoggle{showdiagrams}
\toggletrue{showtoc}
\togglefalse{showtables}
\toggletrue{showfigures}
\togglefalse{showoszillos}
\togglefalse{showdiagrams}

% --- Globale Schriftart für den Haupttext aktivieren ---
\AtBeginDocument{\haupttextSchrift}
% =============================================================================
% EIGENE BEFEHLE & GLOBALE EINSTELLUNGEN (befehle.tex)
% -----------------------------------------------------------------------------
% Dies ist deine zentrale Steuerdatei für das gesamte Dokument.
% Hier passt du globale Variablen, Schriftarten, Abstände und
% Standardgrößen an.
% =============================================================================

% --- GLOBALE DEFINITIONEN FÜR DIE TITELSEITE ---
\newcommand{\praktikumstitel}{Praxisphasenbericht}
\newcommand{\semester}{}
\newcommand{\versuchsnummer}{}
\newcommand{\versuchstitel}{Fremdfeldbeeinflussung auf Messstromwandler in der Niederspannung}
\newcommand{\gruppe}{}
\newcommand{\studiengang}{Studiengang Elektrotechnik}
\newcommand{\vornameStudEins}{}
\newcommand{\nachnameStudEins}{}
\newcommand{\matrikelnummerStudEins}{}
\newcommand{\vornameStudZwei}{Oliver}
\newcommand{\nachnameStudZwei}{Schmidt}
\newcommand{\matrikelnummerStudZwei}{7023462}
\newcommand{\autorenname}{\vornameStudZwei~\nachnameStudZwei}
\newcommand{\betreuerEins}{Dr. Sandro Günter}
\newcommand{\betreuerZwei}{Dipl.-Ing. Holger Kuhlemann}
\newcommand{\betreuerDrei}{Dipl.-Ing. Rainer Ludewig}
\newcommand{\betreuerVier}{Simon Westerbur, B. Eng.}



% --- HINTERGRUNDBILD ---
\newcommand\BackgroundWave{%
    \put(0,0){%
        \parbox[b][\paperheight]{\paperwidth}{%
            \vfill
            \centering
            \vspace{12.0cm}
            \includegraphics[width=\paperwidth]{03_Ressourcen/Logo/hsel-welle-grey}%
            \vfill
        }%
    }%
}


% #############################################################################
% ### ANFANG: GLOBALE LAYOUT- UND SCHRIFTART-EINSTELLUNGEN                  ###
% #############################################################################
%
% Hier kannst du das gesamte Aussehen deines Dokuments anpassen,
% indem du einfach die Werte und Befehle änderst.
% === HAUPTSCHRIFTARTEN ===
% Um die Schriftart zu wechseln, lade das entsprechende Paket in 'praeambel.tex'
% und ändere hier den Befehl.
% Beispiele für Pakete: \usepackage{mathpazo} (Palatino), \usepackage{newtxtext} (Times)
\newcommand{\haupttextSchrift}{\rmfamily}       % \rmfamily (Serif), \sffamily (Sans-Serif)

% === SCHRIFTGRÖSSEN & -STILE ===
% Beispiele: \sffamily (serifenlos), \rmfamily (Serifen), \ttfamily (Schreibmaschine)
%            \bfseries (fett), \itshape (kursiv), \scshape (Kapitälchen)
%            \Huge, \huge, \LARGE, \Large, \large, \normalsize, \small, \footnotesize
\newcommand{\haupttextGroesse}{\normalsize}
\newcommand{\kapitelSchrift}{\sffamily\bfseries\Large}
\newcommand{\sectionSchrift}{\sffamily\bfseries\large}
\newcommand{\subsectionSchrift}{\sffamily\bfseries\normalsize}
\newcommand{\kopfzeilenSchrift}{\haupttextSchrift\small}
\newcommand{\fusszeilenSchrift}{\haupttextSchrift\footnotesize}

% === ABSTÄNDE ===
\newcommand{\zeilenabstand}{1.15}               % 1.0 für einfach, 1.15 für etwas luftiger, 1.5 für 1,5-fachen Abstand
\newcommand{\absatzabstand}{0.5em}              % Abstand zwischen Absätzen (z.B. 1em, 0.5cm, etc.)
\newcommand{\absatzeinzug}{0em}                 % Einzug der ersten Zeile eines Absatzes (0em für linksbündig)

% === KOPF- UND FUSSZEILE ===
\newcommand{\kopfzeilenHoehe}{2.5cm}             % Reservierter Platz für die Kopfzeile
\newcommand{\fusszeilenAbstand}{1.5cm}           % Abstand der Fußzeile vom Text


% =============================================================================
% ### BEFEHLE FÜR AUTOMATISIERTE UMGEBUNGEN ###
% =============================================================================

% Befehl für eine kleine Abbildung (ca. 1/3 der Seitenbreite)
% #1 = Dateipfad
% #2 = Caption (Unterschrift)
% #3 = Label zur Referenzierung
\newcommand{\einfuegenBildKlein}[3]{%
    \begin{figure}[H]
        \centering
        \includegraphics[width=0.4\textwidth]{#1'Bild'}
        \caption{#2}
        \label{#3}
    \end{figure}%
}

% Befehl für eine mittlere Abbildung (ca. 2/3 der Seitenbreite)
% #1 = Dateipfad
% #2 = Caption (Unterschrift)
% #3 = Label zur Referenzierung
\newcommand{\einfuegenBildMittel}[3]{%
    \begin{figure}[H]
        \centering
        \includegraphics[width=0.6\textwidth]{#1}
        \caption{#2}
        \label{#3}
    \end{figure}%
}

% Befehl für eine große Abbildung (fast die ganze Seitenbreite)
% #1 = Dateipfad
% #2 = Caption (Unterschrift)
% #3 = Label zur Referenzierung
\newcommand{\einfuegenBildGross}[3]{%
    \begin{figure}[H]
        \centering
        \includegraphics[width=0.9\textwidth]{#1}
        \caption{#2}
        \label{#3}
    \end{figure}%
}

% #############################################################################
% ### ENDE: GLOBALE LAYOUT- UND SCHRIFTART-EINSTELLUNGEN                    ###
% #############################################################################

% --- LITERATURQUELLE ANGEBEN ---
\addbibresource{bibliographie.bib}

% --- DEFINITIONEN FÜR GLOSSAR LADEN ---
% HINWEIS: Dieser Input-Befehl muss vor \begin{document} stehen!
% =============================================================================
% ABKÜRZUNGSVERZEICHNIS UND GLOSSAR
% -----------------------------------------------------------------------------
% Zentrale Definitionsdatei für alle Abkürzungen und Fachbegriffe.
% Stand: 17.10.2025
% =============================================================================

% --- ABKÜRZUNGEN DEFINIEREN ---

% Phasen
\newacronym{l1}{L1}{Phase L1}
\newacronym{l2}{L2}{Phase L2}
\newacronym{l3}{L3}{Phase L3}

% Organisationen und Normen
\newacronym{vde}{VDE}{Verband der Elektrotechnik, Elektronik und Informationstechnik}

% Technische Begriffe
\newacronym{emv}{EMV}{Elektromagnetische Verträglichkeit}
\newacronym{sps}{SPS}{Speicherprogrammierbare Steuerung}
\newacronym{fem}{FEM}{Finite-Elemente-Methode}
\newacronym{in}{I\textsubscript{N}}{Nennstrom}

% Messverfahren (angepasste, präzisere Definition)
\newacronym{acv}{ACV}{Wechselstrom-Messverfahren (engl. Alternating Current Verfahren)}
\newacronym{dcv}{DCV}{Gleichstrom-Messverfahren (engl. Direct Current Verfahren)}

% --- FACHBEGRIFFE FÜRS GLOSSAR DEFINIEREN ---

\newglossaryentry{sammelschiene}{
    name={Sammelschiene},
    plural={Sammelschienen},
    description={Ein massiver elektrischer Leiter aus Kupfer oder Aluminium in einer Schaltanlage, der zur Verteilung hoher Ströme auf mehrere Abgänge dient.}
}

\newglossaryentry{fremdfeld}{
    name={Fremdfeld},
    plural={Fremdfelder},
    description={Ein magnetisches oder elektrisches Feld, das von einer externen Quelle erzeugt wird und als Störgröße auf ein betrachtetes System einwirkt.}
}

\newglossaryentry{messstromwandler}{
    name={Messstromwandler},
    plural={Messstromwandler},
    description={Ein spezieller Transformator, der hohe Primärströme in ein proportionales, kleines und für Messgeräte verarbeitbares Sekundärsignal umwandelt. Im Fachjargon oft verkürzt als 'Wandler' bezeichnet.}
}

\newglossaryentry{buerde}{
    name={Bürde},
    plural={Bürden},
    description={Die Bürde beschreibt die an den Sekundärkreis eines Messwandlers angeschlossene Last (Impedanz). Sie wird üblicherweise in Voltampere (VA) angegeben und ist entscheidend für die Einhaltung der Genauigkeitsklasse.}
}

\newglossaryentry{genauigkeitsklasse}{
    name={Genauigkeitsklasse},
    plural={Genauigkeitsklassen},
    description={Gibt die maximal zulässige prozentuale Messabweichung eines Messwandlers unter Nennbedingungen an. Gängige Klassen sind z.B. 0,2, 0,5 oder 1.}
}

\newglossaryentry{niederspannung}{
    name={Niederspannung},
    plural={Niederspannungen},
    description={Bezeichnet nach IEC 60038 Wechselspannungen (AC) bis 1000 V und Gleichspannungen (DC) bis 1500 V. In diesem Bereich sind die meisten öffentlichen und industriellen Stromnetze angesiedelt.}
}

\newglossaryentry{induktionsgesetz}{
    name={Induktionsgesetz},
    plural={Induktionsgesetze},
    description={Das Faradaysche Induktionsgesetz ist eine physikalische Grundlage der Elektrotechnik. Es besagt, dass eine zeitliche Änderung des magnetischen Flusses durch eine Leiterschleife eine elektrische Spannung in dieser Schleife induziert. Dies ist das Funktionsprinzip von Transformatoren und somit auch von Messstromwandlern.}
}

\newglossaryentry{niederspannungsschaltanlage}{
    name={Niederspannungsschaltanlage},
    plural={Niederspannungsschaltanlagen},
    description={Eine Anlage zur Verteilung elektrischer Energie im Bereich der Niederspannung. Sie kombiniert Schaltgeräte zum Verbinden und Trennen von Stromkreisen mit Schutz- und Messeinrichtungen wie Sicherungen und Messstromwandlern.}
}

\newglossaryentry{leiter}{
    name={Leiter},
    plural={Leiter},
    description={Ein elektrisch leitendes Material, das dazu dient, elektrischen Strom zu führen. In Schaltanlagen sind dies typischerweise \glspl{sammelschiene} oder Kabel.}
}

\newglossaryentry{erwaermungspruefstand}{
    name={Erwärmungsprüfstand},
    plural={Erwärmungsprüfstände},
    description={Ein spezieller Prüfaufbau zur Verifizierung des thermischen Verhaltens von Schaltanlagen nach \acrshort{vde}-Norm. Mittels eines regelbaren Hochstromtransformators wird der Nenn- oder Überlaststrom bei sehr geringer Spannung (typischerweise unter 10 V) in die Hauptstrombahnen eingespeist. Dies simuliert die reale Wärmeverlustleistung (I²R) im Betrieb und ermöglicht die Überprüfung, ob die entstehende Wärme zuverlässig abgeführt wird und keine unzulässig hohen Temperaturen an Komponenten wie \glspl{sammelschiene} oder Schaltgeräten auftreten.}
}

% --- NEU HINZUGEFÜGTE GLOSSAREINTRÄGE ---

\newglossaryentry{femm}{
    name={FEMM},
    plural={FEMM},
    description={Finite Element Method Magnetics. Eine quelloffene Software zur Simulation von 2D- und achsensymmetrischen Problemen der Elektrostatik und Magnetostatik. Sie wird häufig in der Elektrotechnik zur Analyse von magnetischen Feldern in Geräten wie Transformatoren oder Motoren eingesetzt.}
}

\newglossaryentry{stationaer}{
    name={Stationäres Problem},
    plural={Stationäre Probleme},
    description={Ein Systemzustand, bei dem die betrachteten physikalischen Größen (z.B. Ströme, Felder) zeitlich konstant sind. Dies ist typisch für Gleichstrom-Anwendungen (DC). \gls{femm} ist primär auf die Lösung solcher Probleme ausgelegt.}
}

\newglossaryentry{magnetfeldeinkopplung}{
    name={Magnetfeldeinkopplung},
    plural={Magnetfeldeinkopplungen},
    description={Ein Phänomen, bei dem ein zeitlich veränderliches Magnetfeld einer externen Quelle (z.B. ein benachbarter \gls{leiter}) eine Spannung in einer Leiterschleife (z.B. der Sekundärwicklung eines \gls{messstromwandler}) induziert. Diese Einkopplung ist eine Form der elektromagnetischen Störung und kann die Messgenauigkeit beeinträchtigen.}
}

\newglossaryentry{interne-parameter}{
    name={Interne Parameter},
    plural={Interne Parameter},
    description={Die elektrischen Ersatzgrößen, die das reale Verhalten eines Bauteils wie eines \gls{messstromwandler} beschreiben. Dazu zählen primär der ohmsche Wicklungswiderstand sowie die Haupt- und Streuinduktivitäten.}
}

\newglossaryentry{nennstrom}{
    name={Nennstrom},
    plural={Nennströme},
    description={Der vom Hersteller festgelegte Wert des Primärstroms (\acrshort{in}), für den ein \gls{messstromwandler} ausgelegt ist und bei dem er seine spezifizierte \gls{genauigkeitsklasse} einhält.}
}

\newglossaryentry{genauigkeitsmessung}{
    name={Genauigkeitsmessung},
    plural={Genauigkeitsmessungen},
    description={Ein normiertes Messverfahren zur Überprüfung der Genauigkeit von Messwandlern. Dabei werden der Übersetzungsfehler (Amplitudenfehler) und der Winkelfehler (Phasenfehler) bei definierten Anteilen des \gls{nennstrom} bestimmt und mit den Grenzwerten der \gls{genauigkeitsklasse} verglichen.}
}

\newglossaryentry{acv-verfahren}{
    name={ACV-Messverfahren},
    description={Ein Messverfahren, das Wechselspannung und -strom nutzt, um frequenzabhängige Eigenschaften eines Bauteils zu bestimmen. Bei \glspl{messstromwandler} dient es typischerweise zur Ermittlung von Induktivitäten und Eisenverlusten.}
}

\newglossaryentry{dcv-verfahren}{
    name={DCV-Messverfahren},
    description={Ein Messverfahren, das Gleichspannung und -strom verwendet, um rein ohmsche Widerstände zu messen. Bei einem \gls{messstromwandler} wird es zur präzisen Bestimmung des Kupferwiderstandes der Sekundärwicklung eingesetzt.}
}

% HINZUGEFÜGT, UM DEN FEHLER ZU BEHEBEN
\newglossaryentry{matlab}{
    name={MATLAB},
    description={Eine kommerzielle Software von The MathWorks, Inc. zur Lösung mathematischer Probleme und zur numerischen Simulation. Der Name steht für „Matrix Laboratory“.}
}

% -----------------------------------------------------------------------------
% DOKUMENTENBEGINN
% -----------------------------------------------------------------------------
\begin{document}

% Deckblatt wird als unnummerierte Titelseite eingefügt
% =============================================================================
% DECKBLATT (00_deckblatt.tex)
% -----------------------------------------------------------------------------
% Dies ist die Titelseite deiner Arbeit. Sie verwendet die globalen
% Befehle aus 'befehle.tex', um die Informationen dynamisch einzufügen.
% =============================================================================

\begin{titlepage}

	\hspace{-1.0cm}
	\begin{tabular}{p{8.0cm} p{8.0cm}}
		% KORRIGIERT: Der Pfad zum Logo wurde an die neue Ordnerstruktur angepasst.
		\includegraphics[width = 6.0cm]{03_Ressourcen/Logo/Technik.png} & % Logo
		\parbox[b]{8.0cm}{
		{\large  Fachbereich Technik }                                    \\ % Fachbereich
		{\large  Abteilung Elektrotechnik und Informatik } % Abteilung
		}                                                                 \\
		\\
		\hline
	\end{tabular}
	%
	\begin{center}

		\vspace{2.5cm}
		% KORRIGIERT: Zeilenumbruch \\ entfernt, da \semester leer ist
		\LARGE{\textsc{\praktikumstitel}}
		\ifx\semester\empty\else\LARGE{\textsc{\semester}}\fi

		\vspace{2cm}%
		\ifx\versuchsnummer\empty\else\LARGE{\textsc{\versuchsnummer}}\fi
		\LARGE{\versuchstitel}

		\vspace{2cm}%
		\large

		% KORRIGIERT: Zeilenumbrüche entfernt, da \gruppe leer ist
		\ifx\gruppe\empty\else\gruppe\fi
		\studiengang
		\vspace{0.5cm}

		Vorgelegt von\\

		\begin{table}[!ht]
			\centering
			\begin{tabular}{rll}
				% KORRIGIERT: Leere erste Zeile entfernt
				\vornameStudZwei & \nachnameStudZwei & \matrikelnummerStudZwei % Studierender 2 (global definiert)
			\end{tabular}
		\end{table}


		\vspace{1cm}
		Emden, \today % Ort und Datum

		\vspace{1cm}%
		Betreut von\\ \ \\ % Einleitung Betreuer
		\betreuerEins\\ % Betreuer 1 (global definiert)
		\betreuerZwei\\ % Betreuer 2 (global definiert)
		\betreuerDrei\\ % Betreuer 3 (neu hinzugefügt)
		\betreuerVier

	\end{center}
	\normalsize
\end{titlepage}

% --- VORSAPNN ---
\pagenumbering{Roman} % Römische Seitenzahlen für den Vorspann
% =============================================================================
% ABSTRACT (01_abstract.tex)
% -----------------------------------------------------------------------------
% Eine kurze, prägnante Zusammenfassung deiner gesamten Arbeit.
% (Thema, Vorgehen/Methode, wichtigste Ergebnisse, Schlussfolgerung)
% Sie wird oft auch separat eingereicht.
% =============================================================================

\section*{Abstract}
\addcontentsline{toc}{section}{Abstract}

Hier fasst du deine Arbeit kurz und prägnant zusammen.
% =============================================================================
% VERZEICHNISSE (02_verzeichnisse.tex)
% -----------------------------------------------------------------------------
% Diese Datei generiert alle Verzeichnisse (Inhalt, Abbildungen etc.)
% basierend auf den Schaltern (\toggletrue/\togglefalse) in der Präambel.
% =============================================================================

% Inhaltsverzeichnis
\iftoggle{showtoc}{
    \hypersetup{linkcolor=black}
    \tableofcontents
    \newpage
}{}

% Tabellenverzeichnis
\iftoggle{showtables}{
    \hypersetup{linkcolor=black}
    \listoftables
    \addcontentsline{toc}{section}{Tabellenverzeichnis}
    \newpage
}{}

% Abbildungsverzeichnis
\iftoggle{showfigures}{
    \listoffigures
    \addcontentsline{toc}{section}{Abbildungsverzeichnis}
    \newpage
}{}

% Oszillogrammverzeichnis
\iftoggle{showoszillos}{
    \listofoszillo
    \addcontentsline{toc}{section}{Oszillogrammverzeichnis}
    \newpage
}{}

% Diagrammverzeichnis
\iftoggle{showdiagrams}{
    \listofdiagram
    \addcontentsline{toc}{section}{Diagrammverzeichnis}
    \newpage
}{}

% Setzt die Linkfarbe für den Haupttext wieder auf Blau
\hypersetup{linkcolor=blue}

% --- GLOSSAR & ABKÜRZUNGSVERZEICHNIS AUSGEBEN ---
\printglossary[type=\acronymtype, title={Abkürzungsverzeichnis}]
\printglossary[title={Glossar}]
\newpage

% --- HAUPTTEIL ---
\pagenumbering{arabic} % Arabische Seitenzahlen für den Hauptteil
% =============================================================================
% KAPITEL 1: VORSTELLUNG DES UNTERNEHMENS
% -----------------------------------------------------------------------------
% Hier stellst du das Unternehmen vor, bei dem du die Praxisphase absolviert hast.
% =============================================================================

\section{Vorstellung des Unternehmens}\label{chap:vorstellung_unternehmen}

Das Unternehmen, die Rolf Janssen GmbH, wurde im Jahr 1949 gegründet und blickt auf eine mehr als 75-jährige Geschichte im Bereich der Elektrotechnik zurück. Mit Hauptsitz in Aurich, hat sich das Unternehmen kontinuierlich entwickelt und beschäftigt heute über 200 Mitarbeiter an insgesamt fünf Standorten in Deutschland.

\einfuegenBildGross{03_Ressourcen/Bilder/abbildung_01_firmenzentrale.jpg} {Firmenzentrale der Rolf Janssen GmbH in Aurich} {fig:firmenzentrale}

Das Kerngeschäft der Rolf Janssen GmbH liegt in der Planung, Fertigung und Inbetriebnahme von komplexen elektrotechnischen und leittechnischen Anlagen. Das Leistungsportfolio ist dabei breit diversifiziert und konzentriert sich auf die folgenden Hauptgeschäftsfelder:

\begin{itemize} 
    \item \textbf{Niederspannungsschaltanlagen}: Die Entwicklung und Fertigung von maßgeschneiderten Schaltanlagen bilden eine zentrale Säule des Unternehmens. Hierbei werden individuelle Kundenanforderungen präzise umgesetzt. 
    \item \textbf{Prozessautomatisierung}: Ein weiterer Schwerpunkt ist die Konzeption und Implementierung von modernen Automatisierungslösungen für industrielle Fertigungs- und Verfahrensprozesse. 
    \item \textbf{Diesel-Netzersatzanlagen}: Das Unternehmen stellt Notstromaggregate bereit, die eine unterbrechungsfreie Stromversorgung (USV) für kritische Infrastrukturen und industrielle Anwendungen gewährleisten. 
    \item \textbf{Maritime Lösungen}: Für den anspruchsvollen maritimen Sektor werden spezialisierte Elektro- und Leittechniksysteme entwickelt, die den besonderen Anforderungen auf See (z. B. auf Schiffen oder Offshore-Plattformen) gerecht werden. 
\end{itemize}

Ein wesentliches Alleinstellungsmerkmal der Rolf Janssen GmbH ist die hohe Fertigungstiefe. Diese ermöglicht es, den gesamten Produktionsprozess – vom Rohmaterial bis zum betriebsbereiten Endprodukt – im eigenen Haus zu steuern. Dadurch können Kundenanforderungen nicht nur flexibel, sondern auch kurzfristig umgesetzt werden. Das Unternehmen legt besonderen Wert auf die Realisierung von anspruchsvollen Sonderlösungen und bietet eine durchgängige Projektbetreuung, die von der ersten Angebotsphase über die Konstruktion und Fertigung bis zur finalen Inbetriebnahme reicht. Diese umfassende Herangehensweise positioniert die Rolf Janssen GmbH als kompetenten und zuverlässigen Partner für elektrotechnische Gesamtpakete in diversen Branchen.


% =============================================================================
% KAPITEL 2: PROJEKTBERICHT
% -----------------------------------------------------------------------------
% Das Kernstück deines Berichts. Hier beschreibst du dein Projekt.
% =============================================================================

\section{Bericht über das Projekt}
\label{chap:projektbericht}

\subsection{Motivation}
\label{sec:motivation}
Was war der Anlass für das Projekt? Welches Problem sollte gelöst werden?

\subsection{Vorgehen}
\label{sec:vorgehen}
Wie bist du methodisch vorgegangen? Welche Schritte hast du unternommen?

\subsection{Ergebnisse}
\label{sec:ergebnisse}
Was sind die konkreten Ergebnisse deiner Arbeit?
% =============================================================================
% KAPITEL 3: REFLEXION
% -----------------------------------------------------------------------------
% Hier reflektierst du über deine Erfahrungen während der Praxisphase.
% =============================================================================

\section{Reflexion}
\label{chap:reflexion}

\subsection{Was habe ich gelernt?}
\label{sec:gelernt}
Welche neuen fachlichen und persönlichen Kompetenzen konntest du erwerben?

\subsection{Was lief nicht so gut?}
\label{sec:schwierigkeiten}
Wo gab es Herausforderungen oder unerwartete Probleme?

\subsection{Wie hätte man es besser machen können?}
\label{sec:verbesserung}
Was würdest du im Nachhinein anders machen?
% =============================================================================
% KAPITEL 4: ZUSAMMENFASSUNG UND AUSBLICK
% -----------------------------------------------------------------------------
% Hier fasst du die wichtigsten Punkte zusammen und gibst einen Ausblick.
% =============================================================================

\section{Zusammenfassung und Ausblick}
\label{chap:zusammenfassung_ausblick}

Hier fasst du die zentralen Erkenntnisse deiner Arbeit zusammen und gibst einen Ausblick auf mögliche weiterführende Schritte.


\textcolor{red}{(Damit kein warning angezigt wird \cite{knuth1984texbook})}

% --- ANHANG ---
\appendix
\addsec*{Anhang} % \section{Anhang} ist in der Datei selbst

% --- NACHSPANN ---
\printbibliography[heading=bibintoc, title={Literaturverzeichnis}]
% =============================================================================
% EIGENSTÄNDIGKEITSERKLÄRUNG (eigenstaendigkeitserklaerung.tex)
% -----------------------------------------------------------------------------
% Die formale, unterschriebene Versicherung, dass du die Arbeit
% selbstständig verfasst hast.
% =============================================================================

\section*{Eigenständigkeitserklärung}
\addcontentsline{toc}{section}{Eigenständigkeitserklärung}

Hiermit versichere ich, dass ich die vorliegende Arbeit selbstständig und ohne fremde Hilfe verfasst und keine anderen als die angegebenen Quellen und Hilfsmittel benutzt habe.
\vspace{3cm}

\noindent\makebox[5cm][l]{Ort, Datum} \hspace{4cm} \makebox[5cm][r]{Unterschrift}
\hrule

\end{document}